% Week 08 architecture figure (shop side). Build: xelatex architecture.tex
\documentclass[10pt]{article}
\usepackage[a4paper,landscape,margin=0.9cm]{geometry}
\usepackage{kotex}
\setmainhangulfont{Apple SD Gothic Neo}
\setsanshangulfont{Apple SD Gothic Neo}
\setmainfont{Helvetica Neue}
\usepackage{tikz}
\usetikzlibrary{positioning,arrows.meta,calc,backgrounds,fit}
\pagestyle{empty}

\definecolor{code}{HTML}{2F55A4}
\definecolor{llm}{HTML}{C4601F}
\definecolor{ctx}{HTML}{2F7A5A}
\definecolor{tool}{HTML}{6B46A8}
\definecolor{mkt}{HTML}{6B7280}

\tikzset{
  box/.style={draw=#1, line width=0.8pt, fill=#1!7, rounded corners=2pt, align=left, inner sep=4pt,
              font=\scriptsize, text=black},
  ttl/.style={font=\footnotesize\bfseries},
  band/.style={draw=black!25, fill=black!3, rounded corners=4pt},
  bandlabel/.style={font=\footnotesize\bfseries, text=black!70, anchor=north west},
  lane/.style={font=\scriptsize, text=black!55, anchor=south west},
  arr/.style={-{Stealth[length=2mm]}, line width=0.7pt, draw=black!70},
  rd/.style={-{Stealth[length=2mm]}, line width=0.7pt, draw=black!60, dashed},
  lbl/.style={font=\tiny, text=black!75, fill=white, inner sep=1pt},
  num/.style={circle, draw=black, fill=yellow!60, inner sep=0.6pt, font=\tiny\bfseries},
}
\newcommand{\B}[4]{% name, color, x1 y1, x2 y2 ; content via node text
}

\begin{document}
\centering
\begin{tikzpicture}[x=1cm,y=1cm]

% ---------- title and legend
\node[anchor=west, font=\large\bfseries] at (0,18.6) {Inextro 가게 아키텍처 (Round 3)};
\node[anchor=west, font=\small, text=black!70] at (0,18.0) {LLM은 말하고, 코드가 결정한다. 원본은 전부 남기고, LLM에는 창 · 압축 · 색인 · 교훈만 넣는다.};
\foreach \c/\t [count=\i] in {code/코드(결정), llm/모델(LLM), ctx/컨텍스트(파일), tool/툴, mkt/시장} {
  \node[box=\c, minimum width=1.9cm, align=center] at ({16.3+2.25*\i},18.35) {\t};
}

% ---------- bands
\begin{scope}[on background layer]
  \draw[band] (3.35,2.6) rectangle (12.75,17.4);
  \draw[band] (13.35,2.6) rectangle (19.65,17.4);
  \draw[band] (20.25,2.6) rectangle (27.75,17.4);
\end{scope}
\node[bandlabel] at (3.45,17.35) {빠른 고리: 장사};
\node[bandlabel] at (13.45,17.35) {컨텍스트: 기억과 압축};
\node[bandlabel] at (20.35,17.35) {느린 고리: 매시간 복기 (meta-harness)};
\node[lane] at (3.8,16.65) {매 메시지};
\node[lane] at (3.8,11.25) {매 분 (키트 훅, 60초)};
\node[lane] at (3.8,4.65) {상시};

% ---------- market (left)
\node[box=mkt, minimum width=2.7cm, minimum height=2cm, anchor=south west] (buyer) at (0,13.5)
  {\textbf{buyer 에이전트}\\사람형 · LLM형 ·\\스크립트형};
\node[box=mkt, minimum width=2.7cm, minimum height=2cm, anchor=south west] (kit) at (0,7.9)
  {\textbf{시장 키트}\\재고 · 카탈로그 · 수요\\outbox 전송};
\node[box=mkt, minimum width=2.7cm, minimum height=1.6cm, anchor=south west] (api) at (0,3.0)
  {\textbf{시장 API}\\원장(정산·고장·별점)\\게시판 · 순위};

% ---------- fast loop
\node[box=code, minimum width=4.8cm, minimum height=4.6cm, anchor=south west] (seller) at (3.8,12.0)
  {\textbf{seller.py} (A2A 서버)\\[2pt]
   요청 → SKU · 이름 · 카테고리로 찾기\\
   등급마다 먼저 처분할 단위 1개 (C/B/A)\\
   호가 = 정가 $\times$ 밴딧 팔 (C는 할인)\\
   최저가 = 원가 $\times$ (1+마진) + 보증비용\\
   양보 사다리 → 마지막은 최저가\\
   호가 초과 제안은 1회 확인 후 호가로\\
   체결 → 서명 영수증, claim = 실제 등급\\[2pt]
   \textcolor{black!60}{→ turns · negotiations · missed 로그}};
\node[box=llm, minimum width=3.1cm, minimum height=2.0cm, anchor=south west] (haiku) at (9.2,14.6)
  {\textbf{haiku medium}\\buyer 말 보고 등급 고르기\\답장 쓰기\\\textcolor{llm}{못 봄: 원가 · 최저가 · 정책}};
\node[box=code, minimum width=3.1cm, minimum height=1.4cm, anchor=south west] (check) at (9.2,12.6)
  {\textbf{답장 검사}\\가격 = 끝 숫자, 지어낸 숫자 X\\위반 → 1회 재시도 → template};

\node[box=code, minimum width=8.5cm, minimum height=1.8cm, anchor=south west] (restock) at (3.8,9.4)
  {\textbf{restock.py}\quad 감마-포아송 수요 · 점유율 → newsvendor 임계비로 보유 카테고리 결정\\
   base stock = 다음 몇 분 수요를 임계비 확률로 덮는 양\\
   돈 규칙: 비상금 max(20\%, 100만) · 한 단위 $\le$ 10\% · 비식품 $\le$ 50\%\\
   playbook이 제외한 카테고리는 사지 않음};
\node[box=code, minimum width=8.5cm, minimum height=1.8cm, anchor=south west] (forum) at (3.8,7.2)
  {\textbf{forum.py}\quad 우리 가게 글(이름 · 우리 정산 인용 · 답글)에 모두 답글\\
   \textcolor{llm}{opus medium}이 쓰고, 근거 없는 숫자는 1회 재시도 후 폐기\\
   내부 숫자 대신 buyer의 말로 바꾼 변경만 전달 · 불만은 manager에게};

\node[box=code, minimum width=8.5cm, minimum height=1.6cm, anchor=south west] (coll) at (3.8,3.0)
  {\textbf{collector.py}\quad 원장 · 게시판 · 표 · 잔고 · 순위 · 카탈로그를 계속 수집\\
   해시로 중복 제거, 라운드 내내 전부 보관};

% ---------- context stack
\node[box=ctx, minimum width=5.8cm, minimum height=2.2cm, anchor=south west] (exec) at (13.6,14.4)
  {\textbf{실행}\quad 매시간 덮어씀\\policy.json: 마진 · 양보 · C할인 · 보증 · 돈 규칙\\
   \quad 범위 안, 한 번에 범위의 1/4까지\\playbook.json: 협상 · 게시판 지침, 입고 제외\\
   \quad 숫자 금지 · 600자 · 제외 $\le$ 10개, 24시간};
\node[box=ctx, minimum width=5.8cm, minimum height=2.2cm, anchor=south west] (post) at (13.6,11.8)
  {\textbf{압축}\quad model.py 베이지안 사후분포\\수요(감마-포아송) · 우리 점유율(베타)\\
   등급별 고장률(베타) · 호가 밴딧(톰슨)\\반감기로 오래된 기록의 무게를 줄임};
\node[box=ctx, minimum width=5.8cm, minimum height=3.6cm, anchor=south west] (raw) at (13.6,7.6)
  {\textbf{원본}\quad 전부 보관\\market.db: 모든 가게의 정산 · 고장 · 별점,\\
   \quad 게시판 글과 표, 잔고, 순위\\turns.jsonl: 우리 협상의 모든 말과 답\\
   negotiations · missed · restock 로그};
\node[box=ctx, minimum width=5.8cm, minimum height=2.0cm, anchor=south west] (arch) at (13.6,5.2)
  {\textbf{색인}\quad manager\_archive.jsonl\\시간마다 한 줄: 거래 · 마진 · 놓친 요청 ·\\경고 · 찾은 문제 · 바꾼 것};
\node[box=ctx, minimum width=5.8cm, minimum height=1.8cm, anchor=south west] (mem) at (13.6,3.0)
  {\textbf{교훈}\quad manager\_memory.md (1,500자)\\매시간 다시 씀, 교훈마다 근거 시간};

% ---------- slow loop
\node[box=code, minimum width=6.4cm, minimum height=6.2cm, anchor=south west] (mgr) at (20.8,10.4)
  {\textbf{manager.py} (매시간)\\[2pt]
   1. 브리핑: 숫자 요약 + 작업 창\\
   \quad 지난 1시간 협상 원문 · 놓친 요청\\
   \quad 최근 결정 6개와 각각의 전후 1시간\\
   \quad 프로세스별 마지막 기록 시각 · 경고\\
   \quad 직전 노트 + 교훈 + 색인\\
   2. 필요하면 과거 시간 조회 (최대 2회)\\
   3. 복기와 결정: 지난 결정이 적용되고\\
   \quad 효과가 있었나, 이상은 없나\\
   4. 코드가 걸러서 적용 (clip)\\
   5. 교훈 · 색인 · 사후보고 기록};
\node[box=llm, minimum width=3.0cm, minimum height=2.2cm, anchor=south west] (opus) at (20.8,7.4)
  {\textbf{opus medium}\\(실패 시 sonnet)\\숫자 · 지침 · 교훈\\· 한국어 보고};
\node[box=tool, minimum width=3.0cm, minimum height=2.2cm, anchor=south west] (look) at (24.2,7.4)
  {\textbf{조회 툴}\\lookup(hour, what)\\한 번에 3개 시간대\\DB · 로그 원본 반환};
\node[box=ctx, minimum width=6.4cm, minimum height=1.4cm, anchor=south west] (rep) at (20.8,3.0)
  {\textbf{reports.md}\quad 사후보고 (본 것 · 한 것 · 실제 적용)\\→ 사람이 사후에 읽고 방향을 교정};

% ---------- arrows: market <-> fast loop
\draw[arr, <->] (2.7,14.5) -- node[lbl, above=1pt] {협상} (3.8,14.5);
\draw[arr] (2.0,9.9) |- node[lbl, pos=0.75] {재고} (3.8,12.4);
\draw[arr] (3.8,9.7) -- node[lbl, above=1pt] {구매} (2.7,9.7);
\draw[arr] (3.8,8.3) -- node[lbl, above=1pt] {글} (2.7,8.3);
\draw[arr] (2.7,3.8) -- (3.8,3.8);

% seller <-> haiku <-> check
\draw[arr] (8.6,15.6) -- node[lbl, above=1pt] {사실+수} (9.2,15.6);
\draw[arr] (10.75,14.6) -- (10.75,14.0);
\draw[arr] (9.2,13.2) -- node[lbl, above=1pt] {답장} (8.6,13.2);

% reads: execution + posteriors -> fast loop (bus at x=13.0)
\draw[line width=0.7pt, draw=black!60, dashed] (13.6,15.5) -- (13.0,15.5);
\draw[line width=0.7pt, draw=black!60, dashed] (13.6,12.9) -- (13.0,12.9);
\draw[line width=0.7pt, draw=black!60, dashed] (13.0,15.6) -- (13.0,8.1);
\draw[rd] (13.0,15.6) -- (12.3,15.6);
\draw[rd] (13.0,14.3) -- node[lbl, pos=0.3, above=1pt] {정책 · 지침 · 사후분포 읽기} (8.6,14.3);
\draw[rd] (13.0,10.3) -- (12.3,10.3);
\draw[rd] (13.0,8.1) -- (12.3,8.1);

% writes: logs and collector -> raw
\draw[arr] (12.3,9.8) -- (13.6,9.8);
\draw[arr] (12.3,3.8) -- (13.3,3.8) |- (13.6,8.0);

% slow loop: context -> manager (bus at x=20.05)
\draw[line width=0.7pt, draw=black!70] (20.05,3.9) -- (20.05,11.2);
\draw[arr, <-] (19.4,9.4) -- (20.05,9.4);
\draw[arr, <->] (19.4,6.2) -- (20.05,6.2);
\draw[arr, <->] (19.4,3.9) -- (20.05,3.9);
\draw[arr, <->] (20.05,11.2) -- node[lbl, above=1pt] {} (20.8,11.2);
\draw[arr] (20.8,15.5) -- node[lbl, above=1pt] {갱신} (19.4,15.5);
\draw[arr, <->] (22.3,10.4) -- (22.3,9.6);
\draw[arr] (23.8,8.5) -- node[lbl, above=1pt] {요청} (24.2,8.5);
\draw[arr] (25.7,9.6) -- node[lbl, right=1pt] {원본} (25.7,10.4);
\draw[rd] (25.7,7.4) |- node[lbl, pos=0.75] {원본 조회} (20.05,6.9);
\draw[arr] (27.2,11.0) -- (27.55,11.0) |- (27.2,3.7);

% ---------- numbered differences
\node[num] at (8.45,16.45) {1};
\node[num] at (12.15,11.05) {2};
\node[num] at (27.05,16.45) {3};
\node[num] at (19.25,11.25) {4};

% ---------- how it differs from the kit's starter
\node[anchor=north west, font=\footnotesize\bfseries] at (0,2.35) {키트 기본형과 다른 점};
\node[anchor=north west, text width=13.4cm, font=\scriptsize, align=left] at (0,1.9) {
  \textbf{1} 기본형 가게(seller\_r3)는 원가 이상이면 판다. 여기서는 단위 · 가격 · 체결 · 영수증을 코드가 정하고 haiku는 등급 선택과 문장만 맡는다. 내부 숫자는 LLM 입력에 애초에 없다.\\[2pt]
  \textbf{2} 기본형은 단순한 기본 입고를 쓴다. 여기서는 시장 전체 정산으로 수요를 추정하고(감마-포아송) newsvendor 임계비로 무엇을 얼마나 둘지 정한다.};
\node[anchor=north west, text width=13.4cm, font=\scriptsize, align=left] at (14.2,1.9) {
  \textbf{3} 느린 고리가 하나 더 있다. 매시간 opus가 지난 결정의 효과와 이상을 보고 정책 숫자(범위 안)와 지침(playbook)을 고친다. 코드를 바꾸지 않고 코드가 따르는 값과 프롬프트를 바꾸는 meta-harness다.\\[2pt]
  \textbf{4} 컨텍스트를 층으로 나눴다. 원본은 전부 남기고, 결정에는 사후분포(압축) · 시간별 색인 · 교훈만 넣는다. 필요하면 manager가 조회 툴로 문제 시점의 원본을 꺼내 본다.};

\end{tikzpicture}
\end{document}
